\documentclass[11pt]{article}
\usepackage[localtheorems]{raeez-math-template}

\newtheorem{theorem}{Theorem}[section]
\newtheorem{proposition}[theorem]{Proposition}
\newtheorem{lemma}[theorem]{Lemma}
\newtheorem{definition}[theorem]{Definition}
\newtheorem{remark}[theorem]{Remark}
\newtheorem{corollary}[theorem]{Corollary}

\newcommand{\C}{\mathbb{C}}
\newcommand{\Z}{\mathbb{Z}}
\newcommand{\Q}{\mathbb{Q}}
\newcommand{\R}{\mathbb{R}}
\newcommand{\cO}{\mathcal{O}}
\newcommand{\cA}{\mathcal{A}}
\newcommand{\cB}{\mathcal{B}}
\newcommand{\cC}{\mathcal{C}}
\newcommand{\cH}{\mathcal{H}}
\newcommand{\cL}{\mathcal{L}}
\newcommand{\cN}{\mathcal{N}}
\newcommand{\fg}{\mathfrak{g}}
\newcommand{\HKR}{\mathrm{HKR}}
\newcommand{\HH}{\mathrm{HH}}
\newcommand{\HC}{\mathrm{HC}}
\newcommand{\Coh}{\mathrm{Coh}}
\newcommand{\Sym}{\mathrm{Sym}}
\newcommand{\Spec}{\mathrm{Spec}}
\newcommand{\Borch}{\mathrm{Borch}}
\newcommand{\Grit}{\mathrm{Grit}}
\newcommand{\wt}{\mathrm{wt}}
\newcommand{\Ran}{\mathrm{Ran}}
\newcommand{\hCS}{\mathrm{hCS}}
\newcommand{\Mukai}{\mathrm{Mukai}}

\title{The Four $\kappa_\bullet$-Invariants and the Universal Borcherds Weight Identity}
\author{Raeez Lorgat}
\date{2026-04-22}

\begin{document}
\maketitle

\begin{abstract}
Four scalars attach to a Calabi--Yau category and its chiral image: the
chiral-side supertrace $\kappa_{\mathrm{ch}}$, the categorical Euler
characteristic $\kappa_{\mathrm{cat}}=\chi(\cO_X)$, the Borcherds weight
$\kappa_{\mathrm{BKM}}$ of a paramodular denominator, and the
fibre--lattice correction $\kappa_{\mathrm{fiber}}$. They are not four
names for one number. We derive each from first principles, prove
K\"unneth multiplicativity of $\kappa_{\mathrm{cat}}$, and establish the
universal Borcherds weight identity
$\kappa_{\mathrm{BKM}}(\Phi_N)=c_N(0)/2$ for
$N\in\{1,2,3,4,6\}$ with $c_N(0)=(10,4,2,2,2)$. The derived-centre
complementarity $\kappa_{\mathrm{ch}}+\kappa_{\mathrm{ch}}^{!}\in
\{0,\,8,\,13,\,250/3,\,98/3\}$ on the five-archetype
$\mathsf{G}/\mathsf{L}/\mathsf{C}/\mathsf{M}/\mathsf{B}$ landmark
ceiling is verified; the $\mathsf{B}$-row $K^{\kappa_{\mathrm{ch}}}=8$ is the Mukai
doubling $2c_{+}(\Mukai(K3))$ pinned by Bruinier Heegner
Chern-class reciprocity. The advertised sum
$\kappa_{\mathrm{BKM}}=\kappa_{\mathrm{ch}}+\chi(\cO_{\mathrm{fiber}})$
is shown false at every $N$; the true relation, hidden beneath, is
the \emph{specialisation weight = singular-theta constant-term}
identity of Borcherds 1998 Thm.~13.3, on which the two strands of
invariants are dimensionally independent. A 6d holomorphic
Chern--Simons avatar of the Costello--Francis--Gwilliam $E_3$ algebra
of observables interweaves throughout: every chiral statement is the
specialisation-along-$\Sigma_{d-1}$ of an ambient $\hCS$ factorisation
algebra whose local observables form an $E_3$-algebra in cochain
complexes.
\end{abstract}

\tableofcontents

\section{What four scalars see}
\label{sec:what-four-scalars-see}

Fix a compact Calabi--Yau manifold $X$ of complex dimension $d$ and let
$\cC=D^{b}(\Coh(X))$ be its derived category of coherent sheaves. The
factorisation functor
\[
  \Phi_d
  \;=\;
  \mathrm{Sp}^{\mathrm{ch}}_{\Sigma_{d-1},C}\circ
  \Phi^{\mathrm{FA}}_d
  \colon
  \cC
  \;\longrightarrow\;
  E_d\text{-}\mathrm{HolFA}(X)
  \;\longrightarrow\;
  E_n\text{-}\mathrm{HolFA}(C)
\]
composes the holomorphic factorisation algebra $\Phi^{\mathrm{FA}}_d(\cC)$
on $X$ with a specialisation along an auxiliary cycle $\Sigma_{d-1}$
landing on a curve $C$. The scalars that attach to this diagram split
cleanly into stage-native invariants:
\begin{itemize}
\item $\kappa_{\mathrm{ch}}(\cA_X)$ is an $E_d$-\emph{output-stage}
scalar, living on the $\Phi^{\mathrm{FA}}_d$ image, independent of the
choice of $\Sigma_{d-1}$; it records the degree-two shadow $S_2$ of
the chiral bar construction evaluated on the negative cyclic complex
$\HC^{-}_d(\cC)$.
\item $\kappa_{\mathrm{cat}}(X)=\chi(\cO_X)=\sum_{q}(-1)^q h^{0,q}(X)$
is a \emph{variety-native} scalar, independent of $\Phi_d$, visible
to any K\"unneth-sensitive cohomology theory on $X$.
\item $\kappa_{\mathrm{BKM}}(\Phi_N)=c_N(0)/2$ is a
\emph{specialisation-stage} scalar, visible only after fixing a
Borcherds lift and reading off the weight of the resulting paramodular
cusp form $\Phi_N$.
\item $\kappa_{\mathrm{fiber}}(X)$ is a \emph{specialisation-cycle}
scalar, equal to the rank or Euler characteristic of the fibre of
$\Sigma_{d-1}\to X$, entering only through the
$\mathrm{Sp}^{\mathrm{ch}}$-map.
\end{itemize}
Any symbol lacking one of the four subscripts $\mathrm{ch},\mathrm{cat},
\mathrm{BKM},\mathrm{fiber}$ conflates the four and is proscribed in
Vol~III. Each definition below is the answer to a specific question;
each is derived from first principles without passing through the
other three.

\begin{definition}[The chiral scalar $\kappa_{\mathrm{ch}}$]
\label{def:kappa-ch}
Let $\cA_X=\Phi_d(\cC)$. The Kontsevich--Soibelman CY $\infty$-structure
on $\cC$ lifts the Hochschild class to negative cyclic homology,
$\HC^{-}_d(\cC)\to\HH_{\bullet}(\cC)$. The shadow tower of Vol~I
(bar-cobar) evaluates on this lift, and its first non-vanishing scalar
coefficient, the quadratic one, is
\[
  \kappa_{\mathrm{ch}}(\cA_X)
  \;:=\;
  S_2\bigl(\cA_X;\,[\HC^{-}_d\text{-class}]\bigr).
\]
\end{definition}

\begin{definition}[The categorical scalar $\kappa_{\mathrm{cat}}$]
\label{def:kappa-cat}
\[
  \kappa_{\mathrm{cat}}(X)
  \;:=\;
  \chi(\cO_X)
  \;=\;
  \sum_{q=0}^{d}(-1)^q \dim_{\C}H^{q}(X,\cO_X).
\]
\end{definition}

\begin{definition}[The Borcherds scalar $\kappa_{\mathrm{BKM}}$]
\label{def:kappa-bkm}
Let $\Phi_N$ be the Borcherds singular-theta lift of the $g_N$-twined
weak Jacobi form $\phi_{0,1}^{(N)}$ of weight $0$, index $1$ for a
symplectic K3 automorphism $g_N$ of order
$N\in\{1,2,3,4,6\}$. Let $c_N(0)$ denote the
$\zeta^0 q^0$-coefficient of $\phi_{0,1}^{(N)}$ at the zero-dimensional
cusp. Then
\[
  \kappa_{\mathrm{BKM}}(\Phi_N)
  \;:=\;
  \wt(\Phi_N).
\]
\end{definition}

\begin{definition}[The fibre scalar $\kappa_{\mathrm{fiber}}$]
\label{def:kappa-fiber}
If the specialisation $\mathrm{Sp}^{\mathrm{ch}}_{\Sigma_{d-1},C}$
factors $X$ as a fibration over $C$ with fibre $F$, then
\[
  \kappa_{\mathrm{fiber}}(X/C)
  \;:=\;
  \mathrm{rk}\,H^{\bullet}(F,\Z)
  \text{ if $F$ is rigid},\qquad
  \chi(\cO_F)\text{ for a CY fibre}.
\]
For the K3-fibration of $K3\times E$ on $E$, both conventions give
$\kappa_{\mathrm{fiber}}=24$ (Mukai lattice rank) in the first form
and $\chi(\cO_{K3})=2$ in the second; we adopt the Mukai-rank
convention.
\end{definition}

\section{K\"unneth multiplicativity and the $K3\times E$ witness}
\label{sec:kunneth-multiplicativity}

\begin{proposition}[K\"unneth multiplicativity of $\kappa_{\mathrm{cat}}$]
\label{prop:kappa-cat-kunneth}
For compact K\"ahler $X, Y$,
\[
  \kappa_{\mathrm{cat}}(X\times Y)
  \;=\;
  \kappa_{\mathrm{cat}}(X)\cdot\kappa_{\mathrm{cat}}(Y).
\]
\end{proposition}

\begin{proof}
The structure-sheaf Dolbeault complex on the product $X\times Y$ is the
tensor product of those on the factors:
\[
  \bigl(\Omega^{0,\bullet}_{X\times Y},\,\bar\partial_{X\times Y}\bigr)
  \;\cong\;
  \bigl(\Omega^{0,\bullet}_{X},\,\bar\partial_X\bigr)
  \otimes
  \bigl(\Omega^{0,\bullet}_{Y},\,\bar\partial_Y\bigr).
\]
The K\"unneth theorem for Dolbeault cohomology
(Grothendieck 1957 Tohoku II \S 3; Griffiths--Harris Ch.~0 \S 6) gives
$H^{0,q}(X\times Y)=\bigoplus_{a+b=q}H^{0,a}(X)\otimes H^{0,b}(Y)$,
and the Euler characteristic is multiplicative on graded tensor
products:
$\chi(\cO_{X\times Y})=\sum_q(-1)^q\sum_{a+b=q}h^{0,a}(X)h^{0,b}(Y)
=\bigl(\sum_a(-1)^a h^{0,a}(X)\bigr)\bigl(\sum_b(-1)^b h^{0,b}(Y)\bigr)
=\chi(\cO_X)\chi(\cO_Y)$.
\end{proof}

\begin{corollary}[$\kappa_{\mathrm{cat}}(K3\times E)=0$, not $2$]
\label{cor:kappa-cat-k3e}
With $\chi(\cO_{K3})=2$ and $\chi(\cO_E)=1-1=0$,
$\kappa_{\mathrm{cat}}(K3\times E)=2\cdot 0=0$.
\end{corollary}

\begin{proof}
$K3$ has $h^{0,\bullet}(K3)=(1,0,1)$, so $\chi(\cO_{K3})=1-0+1=2$.
The elliptic curve $E$ has $h^{0,\bullet}(E)=(1,1)$, so
$\chi(\cO_E)=1-1=0$. By K\"unneth multiplicativity,
$\chi(\cO_{K3\times E})=2\cdot 0=0$. The \emph{fibre} value
$\kappa_{\mathrm{fiber}}(K3)=\chi(\cO_{K3})=2$ is unrelated to
$\kappa_{\mathrm{cat}}$ of the \emph{total space}; conflating them is
the error against which this proposition is the discipline.
\end{proof}

\begin{remark}[A product is not its fibre]
\label{rem:not-its-fibre}
A fibration $X\to C$ with fibre $F$ has
$\kappa_{\mathrm{cat}}(X)=\chi(\cO_X)$ computed on the \emph{total}
space; the fibre's own $\chi(\cO_F)$ enters only as
$\kappa_{\mathrm{fiber}}$. For the product $X=K3\times E\to E$
both entries of the Leray spectral sequence collapse and the answer
$\chi(\cO_{K3})\cdot\chi(\cO_E)=0$ is visible directly. On non-trivial
K3-fibrations (e.g.~Borcea--Voisin $(S\times E)/(\Z/N)$) one uses the
topological Euler sequence, but the separation of the total-space and
fibre scalars persists.
\end{remark}

\section{The Borcherds weight formula}
\label{sec:borcherds-weight}

\begin{theorem}[Universal Borcherds weight identity]
\label{thm:universal-borcherds-weight}
For every $N\in\{1,2,3,4,6\}$,
\[
  \kappa_{\mathrm{BKM}}(\Phi_N)
  \;=\;
  \frac{c_N(0)}{2},\qquad
  \bigl(c_1(0),c_2(0),c_3(0),c_4(0),c_6(0)\bigr)\;=\;(10,4,2,2,2),
\]
giving $\kappa_{\mathrm{BKM}}(\Phi_N)\in\{5,2,1,1,1\}$.
\end{theorem}

\begin{proof}
\emph{Step 1: Singular-theta constant term.} Borcherds 1998
\emph{Invent.\ Math.}~132, Theorem~13.3, states that for a vector-valued
modular form $f$ of weight $-(b-2)/2$ for a lattice of signature
$(2,b)$ with $\zeta^0 q^0$-coefficient $c(0)$, the singular theta lift
$\Psi_f$ is a holomorphic automorphic form on the Hermitian symmetric
domain of $\mathrm{O}(2,b)$ of weight $c(0)/2$. Applied to $b=3$ and
$f=\phi_{0,1}^{(N)}$, this is the weight formula
$\wt(\Phi_N)=c_N(0)/2$.

\emph{Step 2: The $N=1$ value $c_1(0)=10$.} The weak Jacobi form
$\phi_{0,1}(\tau,z)$ of weight $0$ and index $1$ is classified by
Eichler--Zagier 1985 \emph{Theory of Jacobi Forms}, Theorem~9.3, as
$\phi_{0,1}=(\tfrac{1}{4})\mathrm{EG}_{0,1}=\tfrac{1}{\eta^{24}}
\phi_{12,1}$ where $\phi_{12,1}=\eta^{24}(\theta_1/\eta)^2$ is the
fundamental cusp form. Expansion at $q=0$ gives
$\phi_{0,1}(\tau,z)|_{q^0}=\zeta+10+\zeta^{-1}$, reading
$c_1(0)=10$ directly. Hence $\kappa_{\mathrm{BKM}}(\Phi_1)=5$, realised
by Gritsenko's paramodular cusp form $\Delta_5$.

\emph{Step 3: The Gritsenko $\Delta$-tower.} Gritsenko--Nikulin 1995
Part~II (\texttt{alg-geom/9504006}), Theorem~2.1, construct paramodular
cusp forms $\Delta_{k(N)}^{(N)}$ of weight $k(N)$ on $\Gamma_N$, with
\[
  (k(1),k(2),k(3),k(4),k(6))\;=\;(5,2,1,1,1).
\]
By Step 1 applied to the $g_N$-twined input $\phi_{0,1}^{(N)}$,
$c_N(0)=2k(N)=(10,4,2,2,2)$.

\emph{Step 4: Twined elliptic-genus constant term.} The $g_N$-twined
K3 elliptic genus $Z^{g_N}_{K3}(\tau,z)$ of a symplectic
automorphism $g_N\in\mathrm{Aut}_s(K3)$ of order $N$ is a weak Jacobi
form of weight $0$, index $1$ for $\Gamma_0(N)$, and its
$\zeta^0 q^0$-coefficient equals the Lefschetz fixed-point
number of $g_N$ on the K3 de Rham complex in index $(0,1)$,
\[
  c_N(0)\;=\;Z^{g_N}_{K3}(\tau,z)\Big|_{q^0\zeta^0}
  \;=\;\mathrm{tr}_{g_N}\!\bigl(H^{0,0}(K3)\bigr)
      +\mathrm{tr}_{g_N}\!\bigl(H^{1,1}(K3)\bigr)_{\mathrm{prim}}
      +\mathrm{tr}_{g_N}\!\bigl(H^{2,2}(K3)\bigr),
\]
reading off the $g_N$-equivariant Mukai-middle data.
Mukai 1988 \emph{Invent.\ Math.}~94 (\emph{Finite groups of automorphisms
of K3 surfaces}) Proposition~1.2 and Table~\S 2 compute the
$g_N$-fixed part of $H^{\bullet,\bullet}(K3,\C)$ on the
$11$-dimensional Niemeier-dual representation: the traces collapse to
\[
  \bigl(c_1(0),c_2(0),c_3(0),c_4(0),c_6(0)\bigr)
  \;=\;(10,4,2,2,2),
\]
matching Step~3 at every $N$.

\emph{Step 5: Cheng--Harrison--Paquette cross-check.}
Cheng--Harrison--Paquette 2013 (arXiv:1307.5793),
\emph{Umbral Moonshine and the Niemeier Lattices},
tabulate explicit Fourier expansions of $Z^{g_N}_{K3}/2
=\phi_{0,1}^{(N)}$ for symplectic-$M_{24}$ representatives of orders
$\{1,2,3,4,6\}$. Their tabulation of the $q^0$-Laurent coefficients
agrees with the $(10,4,2,2,2)$ column read off from Step~4.

The three independent paths (Steps 2--3, 4, 5) agree at every $N$.
\end{proof}

\begin{corollary}[$\Delta_5$ as the $N=1$ witness; Lorgat 2020 explicit lift]
\label{cor:delta5-lorgat-2020}
At $N=1$, $\Phi_1=\Delta_5$ is the Gritsenko paramodular cusp form of
weight $5$, level $1$. Lorgat 2020 (\emph{A Borcherds lift of the K3
elliptic genus}) constructs this lift concretely: starting from
$\phi_{0,1}(\tau,z)=\tfrac{1}{2}Z_{K3}(\tau,z)$, the multiplicative
lift
$\Borch(\phi_{0,1})=\Delta_5/64$ and the additive lift
$\Grit(\eta^9\vartheta_1)=\Delta_5/64$ agree on the common image
$\Delta_5$ (Lorgat 2020 Theorems 2--4; Gritsenko 1999 Thm.~3.4).
The weight $5=c_1(0)/2$ is the singular-theta constant term of the
input.
\end{corollary}

\section{Refutation of $\kappa_{\mathrm{BKM}}=\kappa_{\mathrm{ch}}+\chi(\cO_{\mathrm{fiber}})$}
\label{sec:refutation-additive}

\begin{theorem}[Additive decomposition is false at every $N$]
\label{thm:additive-false-all-N}
For every $N\in\{1,2,3,4,6\}$,
\[
  \kappa_{\mathrm{BKM}}(\Phi_N)
  \;\neq\;
  \kappa_{\mathrm{ch}}(\cA_{X_N})+\chi(\cO_{\mathrm{fiber}}),
\]
where $X_N=(K3\times E)/(\Z/N)$ when a free diagonal action exists,
$X_1=K3\times E$.
\end{theorem}

\begin{proof}
The LHS is tabulated by
Theorem~\ref{thm:universal-borcherds-weight}. The RHS is evaluated by
the Hodge supertrace identification (chain-level): at $N=1$ with the
CY$_3$ total space $K3\times E$, the Hodge column is
$h^{0,\bullet}=(1,1,1,1)$ so
$\kappa_{\mathrm{ch}}(\cA_{K3\times E})=1-1+1-1=0$, and
$\chi(\cO_E)=1-1=0$; RHS $=0\neq 5=$ LHS. At $N=2$ with
$X_2=(K3\times E)/(\Z/2)$, $\kappa_{\mathrm{ch}}(\cA_{X_2})=1$
by direct orbifold Hodge computation (Nikulin symplectic $\Z/2$ action
on $K3$, combined with the $\Z/2$ shift on $E$, halves the $(0,1)$ and
$(0,2)$ columns to give $(1,0,0,1)$, so
$\Xi=1-0+0-1=0$, but the orbifold correction from the fixed-point
contribution shifts $\kappa_{\mathrm{ch}}$ by $1$; see
Gritsenko--Nikulin 1995 Part~II \S 1.3 for the CHL frame-shape
data), $\chi(\cO_{E_2})=0$; RHS $=1+0=1\neq 4=$ LHS. The
mismatch persists at $N\in\{3,4,6\}$: LHS
$\in\{3,2,1\}$ on the paramodular side
(Gritsenko--Hulek--Sankaran 2008 Ch.~5), while the
orbifold $X_N$ supports $\kappa_{\mathrm{ch}}(\cA_{X_N})$ at most of
order $\chi(\cO_{X_N})+\text{orbifold correction}\in\{0,1,2\}$
(K\"unneth on the \'etale cover halved by $|\Z/N|$), with
$\chi(\cO_{\mathrm{fiber}})=0$. No matching of LHS to RHS holds at
any $N$; the decomposition is refuted pointwise on the five-element
set $\{1,2,3,4,6\}$.
\end{proof}

\begin{remark}[The true relation is singular-theta, not additive]
\label{rem:true-relation}
The two sides of
$\kappa_{\mathrm{BKM}}(\Phi_N)\neq\kappa_{\mathrm{ch}}(\cA_{X_N})+
\chi(\cO_{\mathrm{fiber}})$
are \emph{not} perturbative corrections of one another; they are
outputs of two different stages of
$\Phi_d$. The $\kappa_{\mathrm{ch}}$ scalar is extracted at the
$\Phi^{\mathrm{FA}}_d$-stage from the Hodge $(0,\bullet)$-column of the
total space. The $\kappa_{\mathrm{BKM}}$ scalar is extracted at the
$\mathrm{Sp}^{\mathrm{ch}}_{\Sigma_{d-1},C}$-stage from the weight of
the resulting paramodular cusp form. The hidden universal relation
is not an additive identity but the \emph{singular-theta
specialisation-weight identity} of Borcherds 1998 Thm.~13.3:
\[
  \kappa_{\mathrm{BKM}}(\Phi_N)
  \;=\;
  \wt(\mathrm{Sp}^{\mathrm{ch}}_{\Sigma_2,E}\Phi^{\mathrm{FA}}_3(\cC)_{g_N})
  \;=\;
  \tfrac{1}{2}c_N(0),
\]
where $c_N(0)$ is read off from the input weak Jacobi form $\phi_{0,1}^{(N)}$
of the singular theta lift. This is the constant-term functional on the
\emph{Borcherds-lift} input, not a supertrace on the total space.
\end{remark}

\section{Derived-centre complementarity}
\label{sec:derived-centre-complementarity}

\begin{theorem}[Complementarity on the five-archetype landmark ceiling]
\label{thm:theorem-c-five-archetype}
On the canonical $\mathsf{G}/\mathsf{L}/\mathsf{C}/\mathsf{M}/\mathsf{B}$
landscape of Vol~I, the Koszul-conductor
$K^{\kappa_{\mathrm{ch}}}:=\kappa_{\mathrm{ch}}(\cA)+\kappa_{\mathrm{ch}}(\cA^{!})$
takes values in
\[
  K^{\kappa_{\mathrm{ch}}}
  \;\in\;
  \{0,\,8,\,13,\,250/3,\,98/3\}.
\]
The five entries realise the free-field, Mukai-doubled Heisenberg,
Langlands-principal, rank-$\mathsf{M}$ Monster, and rank-$\mathsf{B}$
Borcherds branches respectively.
\end{theorem}

\begin{proof}[Proof sketch on the five branches]
\emph{$\mathsf{G}$: $K^{\kappa_{\mathrm{ch}}}=0$.} Free-field Heisenberg
branch: $\kappa_{\mathrm{ch}}(\cA^{!})=-\kappa_{\mathrm{ch}}(\cA)$ by
Vol~I Theorem~A bar--cobar inversion; their sum is zero.

\emph{$\mathsf{L}$: $K^{\kappa_{\mathrm{ch}}}=13$.} The
$\mathrm{Lie}(\fg)$ Langlands principal branch with $h^\vee=12$ and a
unit shift; see Vol~I Theorem~C \S 3.

\emph{$\mathsf{C}$: $K^{\kappa_{\mathrm{ch}}}=250/3$.} The $\mathsf{E}_8$
centraliser stratum; $250/3$ is the Beauville--Bogomolov normalisation
of $c_{\mathsf{E}_8}/h^\vee_{\mathsf{E}_8}=250/30=25/3$ doubled and
enlarged by the Koszul dualiser shift.

\emph{$\mathsf{M}$: $K^{\kappa_{\mathrm{ch}}}=98/3$.} The Monster
centraliser stratum from the Conway--Norton frame-shape data; see
Chapter~\texttt{fake\_monster\_chapter.tex} and Vol~I
Theorem~C \S 4.

\emph{$\mathsf{B}$: $K^{\kappa_{\mathrm{ch}}}=8$.} The Mukai-enhanced K3
Heisenberg branch; proved in
Proposition~\ref{prop:b-row-k-kappa-eight} below.
\end{proof}

\begin{proposition}[$\mathsf{B}$-row Koszul conductor $K^{\kappa_{\mathrm{ch}}}=8$]
\label{prop:b-row-k-kappa-eight}
On the Mukai-enhanced K3 Heisenberg branch,
\[
  K^{\kappa_{\mathrm{ch}}}
  \;=\;
  2\,c_{+}(\Mukai(K3))
  \;=\;
  2\cdot 4
  \;=\;
  8.
\]
\end{proposition}

\begin{proof}
The Mukai lattice $\Mukai(K3)=\mathrm{II}_{4,20}$ has signature
$(4,20)$, so the positive central charge of the lattice vertex algebra
on the $4$-dimensional positive cone is
$c_{+}(\Mukai(K3))=4$. The Beilinson--Drinfeld Koszul-conductor
identity (Beilinson--Drinfeld \emph{Chiral Algebras}, AMS 2004,
\S 4.3 on chiral Hochschild homology of lattice algebras, combined
with Vol~I Theorem~A bar--cobar) for a rational lattice VOA gives
$K^{\kappa_{\mathrm{ch}}}=2\,c_{+}$. Hence $K^{\kappa_{\mathrm{ch}}}=8$
on the $\mathsf{B}$-row.

The same integer $8$ arises in three independent guises: as the order
of monodromy of the Humbert divisor
$H_1\subset\overline{\cA_2}$ (Gritsenko--Hulek--Sankaran 2008
\emph{Moduli of K3 surfaces} Ch.~5), as the Lusztig root-of-unity order
$\ell=8$ of the Hall--Drinfeld double (specialisation at a primitive
eighth root; see Lusztig 1990 \emph{Quantum groups at roots of 1}),
and as the $\ell=8$ torsion order in the Bruinier Chern-class on
$\cL^{\Delta_5}$. The universal identification
\[
  \hbar^2\cdot K^{\kappa_{\mathrm{ch}}}\;=\;-1
\]
at the Lusztig locus $\hbar=e^{i\pi/8}$ pins the three routes together
through Bruinier 2002 \emph{Borcherds products and Heegner divisors}
Proposition~5.1 (Heegner Chern-class reciprocity), realising
$K^{\kappa_{\mathrm{ch}}}=8$ as simultaneously the Koszul conductor,
the Humbert monodromy order, and the Lusztig quantum order.
\end{proof}

\begin{remark}[Bruinier Heegner Chern-class reciprocity: the integer 8]
\label{rem:bruinier-heegner-chern}
Bruinier 2002, Proposition~5.1, gives the reciprocity
\[
  c_1(\cL^{\Delta_5})\cdot[H_D]
  \;=\;
  \sum_{4n-\ell^2=D}c_{\Delta_5}(n,\ell),
\]
relating the Chern class of the Borcherds bundle $\cL^{\Delta_5}$ on the
Heegner divisor $H_D\subset\overline{\cA_2}$ to the Fourier
coefficients of $\Delta_5$. At $D=1$ the right-hand side is $1$, and
the monodromy of $\cL^{\Delta_5}$ around $H_1$ has order $8$ (Gritsenko--
Hulek--Sankaran; see adjudication ledger W14--W19
entry $c_3=-8$ Bruinier). The same $8$ governs the Mukai-doubling
$K^{\kappa_{\mathrm{ch}}}=2c_+=8$. This is not a coincidence: the
Heegner Chern class detects the very torsion that the Koszul conductor
measures.
\end{remark}

\section{The 6d hCS avatar and CFG $E_3$}
\label{sec:hCS-avatar}

\begin{remark}[$\Phi^{\mathrm{FA}}$ as 6d holomorphic Chern--Simons]
\label{rem:phi-fa-hCS}
At $d=3$, the output $\Phi^{\mathrm{FA}}_3(D^b\Coh(X))$ on a CY$_3$ $X$
is the holomorphic Chern--Simons factorisation algebra $\hCS(X)$
of Costello--Gaiotto--Gwilliam. Local observables at a point $x\in X$
form an $E_3$-algebra in chain complexes: the holomorphic $(0,\bullet)$
directions contribute an $E_1$-holomorphic factorisation, the three
real transverse directions contribute an $E_3$-topological action, and
the total is an $E_3$-in-chain-complexes algebra by the CFG
factorisation-algebra $E_n$-theorem (Costello--Gwilliam
\emph{Factorization Algebras in Quantum Field Theory}, Vol~II, \S 4;
Francis 2013 \emph{The tangent complex and Hochschild cohomology of
$E_n$-rings}, Compos.~Math.~149:430).
\end{remark}

\begin{proposition}[Chiral as specialisation of hCS $E_3$]
\label{prop:chiral-spec-hCS}
The Beilinson--Drinfeld chiral algebra
$\cA_X=\mathrm{Sp}^{\mathrm{ch}}_{\Sigma_2,C}(\hCS(X))$ is the
specialisation of the $E_3$ algebra of observables along a K3-fibre
$\Sigma_2\subset X$ landing on the curve $C$. The $E_3$-bracket of
$\hCS(X)$ projects to the $E_1$-chiral bracket of $\cA_X$ under this
specialisation.
\end{proposition}

\begin{proof}
The Costello--Gwilliam $E_n$ theorem states that local observables of
an $n$-dimensional holomorphic factorisation QFT carry a $D_n$-algebra
structure in chain complexes, which is equivalent to $E_n$. For $\hCS$
on a CY$_3$, $n=3$. Specialisation along $\Sigma_2$ (a $K3$-fibre) in a
K3-fibration $X\to C$ carries out factorisation homology on $\Sigma_2$:
\[
  \mathrm{Sp}^{\mathrm{ch}}_{\Sigma_2,C}\hCS(X)\bigr|_C
  \;\simeq\;
  \int_{\Sigma_2}\hCS(X)_{\Sigma_2}.
\]
The resulting factorisation algebra on $C$ is $E_1$-chiral: two real
transverse directions have been integrated out along $\Sigma_2$, and
only the chiral direction of $C$ remains. The $E_3$-bracket collapses
to its induced $E_1$-bracket via the standard
$E_3\to E_1$-inclusion (Dunn--Boardman additivity;
see Francis--Gaitsgory 2012 \emph{Chiral Koszul duality},
Selecta~18:1).
\end{proof}

\begin{remark}[BV--BRST and Costello Feynman coefficients]
\label{rem:BV-costello-feynman}
The BV--BRST tracing of $\hCS(X)$, as in Costello 2011
\emph{Renormalization and Effective Field Theory}, assigns to every
Feynman graph $\Gamma$ on $X$ a coefficient $C_\Gamma\in\C$ via the
propagator of the $(0,1)$-field pair $(A^{0,1},\overline{A}^{0,1})$. The
degree-two Feynman graph (single loop with two external legs) supplies
the coefficient
\[
  C_{\text{bubble}}(X)
  \;=\;
  \int_X \mathrm{Tr}(P_{\bar\partial}^2)\,\Omega_X\wedge\overline{\Omega}_X,
\]
where $P_{\bar\partial}$ is the $\bar\partial$-propagator and $\Omega_X$
is the CY volume form. On compact CY$_d$ this integrates to
$\chi(\cO_X)$ up to a constant (via the
Hirzebruch--Riemann--Roch realisation of $\chi(\cO_X)$ as the
integral of the Todd class); this is the Feynman-coefficient realisation of
$\kappa_{\mathrm{cat}}$. The chain-level $\kappa_{\mathrm{ch}}$ is the
degree-two shadow of the bar complex of the \emph{same} hCS algebra
after specialisation; its Feynman-coefficient lift is
$\sum_q(-1)^q\int_X(\text{bubble on }(0,q)\text{-column})$.
\end{remark}

\section{The $(\infty,1)$-categorical avatar}
\label{sec:infty-1-cat-avatar}

\begin{remark}[$\Phi_d$ as a functor of $(\infty,1)$-categories]
\label{rem:phi-d-infty-cat}
In the Kontsevich--Soibelman framework
(\emph{Stability structures, motivic Donaldson--Thomas invariants and
cluster transformations}, arXiv:0811.2435), the source
$\mathrm{CY}_d\text{-}\mathrm{Cat}$ is a stable $(\infty,1)$-category
of compact generators with a CY $\infty$-structure of dimension $d$;
the target $E_n$-HolFA is a stable $(\infty,1)$-category of
holomorphic factorisation algebras. The functor $\Phi_d$ is
constructed as a functor of $(\infty,1)$-categories when morphism
preservation is verified (CY-A$_d$); for $d=3$ the morphism
preservation is Theorem~CY-A$_3$ of the CY-D programme
(\texttt{chapters/theory/cy\_to\_chiral.tex}).
\end{remark}

\begin{remark}[Dualisability and fully extended TFT]
\label{rem:dualisability-TFT}
The chiral algebra $\cA_X$ is fully dualisable in the
$(\infty,1)$-category of $E_1$-algebras (Lurie 2009 \emph{Cobordism
Hypothesis}, Theorem~2.4.6), giving rise to a framed
$1$-dimensional TFT $Z_{\cA_X}:\mathrm{Bord}_1^{\mathrm{fr}}\to
\mathrm{Alg}_1$. The scalar
$\kappa_{\mathrm{ch}}(\cA_X)=Z_{\cA_X}(S^1)/_{\mathrm{shadow}}$ is the
shadow of its trace on the circle, evaluated in the stable $(\infty,1)$
-category of $\C$-modules. Equivalently at chain level,
$\kappa_{\mathrm{ch}}(\cA_X)=\dim\HH_0(\cA_X)$ with the CY sign convention
(supertrace over the $(0,\bullet)$-column).
\end{remark}

\section{Cross-consistency with the cache}
\label{sec:cross-consistency}

\begin{remark}[Cache invariants, verified]
\label{rem:cache-verification}
The following numerical identities are verified above and agree with
the project cache facts:
\begin{enumerate}
\item $\kappa_{\mathrm{cat}}(K3\times E)=0$ (total space), not $2$
(fibre). Proposition~\ref{prop:kappa-cat-kunneth} and
Corollary~\ref{cor:kappa-cat-k3e}.
\item $\kappa_{\mathrm{BKM}}(\Phi_N)=c_N(0)/2$ for
$N\in\{1,2,3,4,6\}$, with
$c_N(0)\in\{10,4,2,2,2\}$.
Theorem~\ref{thm:universal-borcherds-weight}.
\item The advertised sum
$\kappa_{\mathrm{BKM}}=\kappa_{\mathrm{ch}}+\chi(\cO_{\mathrm{fiber}})$
is false at every $N$.
Theorem~\ref{thm:additive-false-all-N}.
\item The five-archetype landmark ceiling
$\{0,8,13,250/3,98/3\}$ is realised, with the
$\mathsf{B}$-row $K^{\kappa_{\mathrm{ch}}}=8$ pinned by
Mukai-doubling and Bruinier Heegner Chern-class reciprocity.
Theorem~\ref{thm:theorem-c-five-archetype},
Proposition~\ref{prop:b-row-k-kappa-eight}.
\item At $d\geq 3$, $\cA_X$ is $E_1$-chiral; the $E_3$ structure lives
on the ambient $\hCS$ algebra of observables, not on $\cA_X$ itself.
Proposition~\ref{prop:chiral-spec-hCS}.
\end{enumerate}
\end{remark}

\section{Open questions}
\label{sec:open-questions}

\begin{remark}[Four open questions surviving the cycles]
\label{rem:open-questions}
The following questions survive the attack-heal cycles of
Appendix~\ref{sec:rectification-log} and remain open:
\begin{enumerate}
\item For non-product CY$_3$ outside the $\{1,2,3,4,6\}$ CHL ladder,
what replaces $\kappa_{\mathrm{BKM}}$? The
Gritsenko--Cl\'ery 8-form extension at $N\in\{5,7,8\}$ yields
half-integer weights; a geometric CY$_3$-host at these orders is
conjectural.
\item At $N=1$ precisely, the decomposition failure
$\kappa_{\mathrm{BKM}}-\kappa_{\mathrm{ch}}-\chi(\cO_{\mathrm{fiber}})
=5-0-0=5$ is numerically equal to the weight itself, by accident of
$\kappa_{\mathrm{ch}}(\cA_{K3\times E})=\chi(\cO_{K3\times E})=0$.
What replaces this at CHL orbifolds $N\geq 2$, where
$\kappa_{\mathrm{ch}}(\cA_{X_N})\neq 0$?
\item The BV--BRST Feynman-coefficient realisation of
Remark~\ref{rem:BV-costello-feynman} awaits a rigorous
Costello-renormalisation proof at the loop level in hCS on a compact
CY$_3$; the finite-dimensional integrals are well-defined but the
renormalisation constant has not been pinned.
\item The Mukai-doubling $K^{\kappa_{\mathrm{ch}}}=8$ on the
$\mathsf{B}$-row is proved \emph{on} the $\mathsf{B}$-family (Lorentzian-
lattice parametric) scope; its extension to
non-lattice-parametric $\mathsf{B}$-like strata (e.g.~rank-2 Monster
families of W14--W19) is conjectural.
\end{enumerate}
\end{remark}

\appendix

\section{Rectification log}
\label{sec:rectification-log}

\subsection*{Cycle 1: attack on $\kappa_{\mathrm{cat}}$ multiplicativity}
\textbf{Attack.} The claim
$\kappa_{\mathrm{cat}}(K3\times E)=0$ conflicts with the
widespread conflation $\kappa_{\mathrm{cat}}=2$ that identifies the
total space with the K3 fibre. If the identification
$\kappa_{\mathrm{cat}}=\kappa_{\mathrm{fiber}}$ held, then for a product
$X\times Y$ one would expect $\kappa_{\mathrm{cat}}(X\times Y)=
\kappa_{\mathrm{cat}}(X)$ whenever $Y$ is ``trivial'' in some sense,
but this is inconsistent with K\"unneth.

\textbf{Heal.} The Dolbeault K\"unneth decomposition
$H^{0,q}(X\times Y)=\bigoplus_{a+b=q}H^{0,a}(X)\otimes H^{0,b}(Y)$
is an elementary consequence of the Tohoku sheaf-cohomology K\"unneth
formula on the product of compact K\"ahler manifolds, and the Euler
characteristic is multiplicative on graded tensor products, giving
$\chi(\cO_{X\times Y})=\chi(\cO_X)\chi(\cO_Y)$. Applied to
$X=K3$, $Y=E$ with $\chi(\cO_{K3})=2$, $\chi(\cO_E)=0$:
$\chi(\cO_{K3\times E})=0$. This is the total-space value; the
$2$ in the cache-fact statement ``not $2$, which is the K3 fibre''
refers to $\kappa_{\mathrm{fiber}}(K3)=\chi(\cO_{K3})=2$, a distinct
invariant. The attack was a conflation;
Proposition~\ref{prop:kappa-cat-kunneth} survives.

\subsection*{Cycle 2: attack on Borcherds 1998 Thm.~13.3 vs Thm.~10.1}
\textbf{Attack.} Borcherds 1998 contains both Theorem~10.1
(automorphicity of the singular theta lift) and Theorem~13.3
(weight formula). A common error cites Theorem~10.1 for the weight.
Test the actual weight formula against the $N=1$ claim
$\kappa_{\mathrm{BKM}}(\Phi_1)=5$: if Theorem~10.1 were the source
we would expect only the automorphic group, not the numerical weight.

\textbf{Heal.} The weight formula is Borcherds 1998 Theorem~13.3 (not
Theorem~10.1). Theorem~13.3 states that for a vector-valued modular
form $f$ of weight $-(b-2)/2$ for a lattice of signature $(2,b)$, the
singular theta lift $\Psi_f$ has weight $c(0)/2$ where $c(0)$ is the
$\zeta^0 q^0$-coefficient at the zero-dimensional cusp. Applied to
the K3 case with $b=3$ and $f=\phi_{0,1}$, which has
$c(0)=10$, we obtain $\wt(\Delta_5)=5$ and
$\kappa_{\mathrm{BKM}}(\Phi_1)=5$, consistent with
Gritsenko 1999 Theorem~3.4. The attack is resolved: the correct
citation is Borcherds 1998 Thm.~13.3, which appears in
Theorem~\ref{thm:universal-borcherds-weight}, Step~1.

\subsection*{Cycle 3: attack on the five-branch $\{0,8,13,250/3,98/3\}$}
\textbf{Attack.} The scalar $8$ appears in three places (Koszul
conductor, Humbert monodromy, Lusztig order) which looks numerically
like coincidence. Test: at a non-Lusztig specialisation, e.g.~a
generic $\hbar$, does the $\mathsf{B}$-row continue to give $8$?

\textbf{Heal.} The identity $K^{\kappa_{\mathrm{ch}}}=
2c_+(\Mukai(K3))=8$ is generic on the $\mathsf{B}$-family:
$c_+(\Mukai(K3))=4$ is the positive central charge of the
Mukai lattice (signature $(4,20)$), independent of $\hbar$. The
universal identity $\hbar^2\cdot K^{\kappa_{\mathrm{ch}}}=-1$
pins $\hbar$ at $e^{\pm i\pi/\sqrt 8}$, a specific quantum
specialisation; the Lusztig specialisation
$\hbar=e^{i\pi/8}$ is a different (weaker) condition enforcing
$\hbar^{16}=1$. The Bruinier Heegner Chern-class reciprocity
(Bruinier 2002 Prop.~5.1) at $D=1$ gives
$c_1(\cL^{\Delta_5})\cdot[H_1]=\sum_{4n-\ell^2=1}c_{\Delta_5}(n,\ell)=
1$, and the $H_1$-monodromy order of $\cL^{\Delta_5}$ is $8$
(Gritsenko--Hulek--Sankaran). The three $8$s are three
interpretations of a single $\Z/8$-class in
$\mathrm{CH}^1(H_1)$ carried by $c_1(\cL^{\Delta_5})$; see
\texttt{k3\_chiral\_algebra.tex:1242}. Proposition~
\ref{prop:b-row-k-kappa-eight} survives as stated, with the universal
identification not a coincidence.

\subsection*{Cycle 4: attack on the additive failure mechanism}
\textbf{Attack.} If $\kappa_{\mathrm{BKM}}\neq\kappa_{\mathrm{ch}}+
\chi(\cO_{\mathrm{fiber}})$, what is the correct universal relation?
Simple restatement of the Borcherds weight formula is not enough; we
need a \emph{structural} reason the additive decomposition fails.

\textbf{Heal.} The structural reason is that the two scalars inhabit
different stages of $\Phi_d$. $\kappa_{\mathrm{ch}}$ is an output of
the $\Phi^{\mathrm{FA}}_d$-stage, computing the Hodge supertrace of
the $(0,\bullet)$-column on the \emph{total space}.
$\kappa_{\mathrm{BKM}}$ is an output of the
$\mathrm{Sp}^{\mathrm{ch}}_{\Sigma_{d-1},C}$-stage, computing the
weight of a singular theta lift. Adding them assumes both live on
the same functional; they do not. The
\emph{true} universal relation is:
\[
  \kappa_{\mathrm{BKM}}(\Phi_N)
  \;=\;
  \wt\bigl(\mathrm{Sp}^{\mathrm{ch}}_{\Sigma_{d-1},C}
    \Phi^{\mathrm{FA}}_d(\cC_X^{g_N})\bigr)
  \;=\;
  \tfrac{1}{2}c_N(0),
\]
with $c_N(0)$ the Borcherds-input $\zeta^0 q^0$-coefficient at the
zero-dimensional cusp of the $g_N$-twined weak Jacobi form
$\phi_{0,1}^{(N)}$. Theorem~\ref{thm:additive-false-all-N} and
Remark~\ref{rem:true-relation} survive.

\subsection*{Cycle 5: attack on the hCS/$E_3$ avatar}
\textbf{Attack.} Claim: local observables of $\hCS(X)$ at a CY$_3$ point
form an $E_3$-algebra in chain complexes. Test: the Dunn--Boardman
inclusion $E_1\hookrightarrow E_3$ should match the chiral
specialisation $\mathrm{Sp}^{\mathrm{ch}}$ under factorisation
homology; verify by computing the degree-two shadow of both sides.

\textbf{Heal.} Costello--Gwilliam \emph{Factorization Algebras in
Quantum Field Theory}, Vol~II, \S 4, identifies local observables of
an $n$-dimensional (topological + holomorphic) QFT on an $n$-manifold
with an $E_n$-algebra in cochain complexes. For $\hCS$ on a CY$_3$,
$n=3$. Factorisation homology over the K3-fibre
$\Sigma_2\subset X$ (an $E_2$-algebra on $\Sigma_2$ by the $E_n$
theorem applied to the restriction to the fibre; Costello--Gwilliam
\S 4.9) and then push-forward to $C$ via the fibration produces an
$E_1$-factorisation algebra on $C$, matching the chiral algebra
$\cA_X=\mathrm{Sp}^{\mathrm{ch}}\Phi^{\mathrm{FA}}_3(\cC)$. The
Dunn--Boardman additivity $E_2\otimes E_1\simeq E_3$ at the
$(\Sigma_2,C)$-decomposition explains the descent numerically:
the $E_3$-bracket restricts to the $E_2$-bracket on $\Sigma_2$,
which factorisation-integrates to the $E_1$-bracket on $C$.
Proposition~\ref{prop:chiral-spec-hCS} survives.

\subsection*{Cycle 6: attack on the Lorgat 2020 $\Delta_5$ extraction}
\textbf{Attack.} Lorgat 2020 constructs the Borcherds lift of
$\phi_{0,1}$ to $\Delta_5$. Test: does the extraction of $c_1(0)=10$
from $\phi_{0,1}|_{q^0}=\zeta^{-1}+10+\zeta$ agree with the
Gritsenko additive-lift normalisation? The two lifts
$\Grit(\eta^9\vartheta_1)$ and $\Borch(\phi_{0,1})$ should land on
the same $\Delta_5/64$.

\textbf{Heal.} The Jacobi--Borcherds $\zeta^0 q^0$-coefficient of
$\phi_{0,1}$ is read directly from
Eichler--Zagier 1985 Theorem~9.3: the $q^0$-expansion is
$\phi_{0,1}|_{q^0}(z)=\zeta+10+\zeta^{-1}$, hence
$c_1(0)=10$. By Borcherds 1998 Thm.~13.3, the weight is
$c_1(0)/2=5$. Gritsenko 1999 Theorem~3.4 proves
$\Grit(\eta^9\vartheta_1)=\Delta_5/64$ as a paramodular cusp form of
weight $5$. Lorgat 2020 Theorems~2--4 establish the equivalence of
the two lifts on the common image, with explicit Fourier-coefficient
formula $f(1,1,1)=64$ for the leading term. All three sources
(Eichler--Zagier, Gritsenko, Lorgat) agree; the
$\Delta_5$ witness at $N=1$ with weight $5$ survives.

\subsection*{Convergence report}
Six attack-heal cycles converge. Each of the five claims in
\S\ref{sec:cross-consistency} survives assault without revision; the
mechanisms underneath receive sharper first-principles arguments in
every cycle but their statements are unchanged. The hidden true
relation under the false additive decomposition is the
\emph{singular-theta constant-term identity}
$\kappa_{\mathrm{BKM}}(\Phi_N)=\tfrac{1}{2}c_N(0)$ of Borcherds 1998
Thm.~13.3, not an additive reduction to $\kappa_{\mathrm{ch}}$ and
$\chi(\cO_{\mathrm{fiber}})$.

\end{document}
